\chapter{Présentation du Projet et Cahier des Charges}
\label{ch:presentation}

\section{Concept général}
\label{sec:presentation_concept}

\esibot{} est un robot mobile autonome à entraînement différentiel (\emph{differential
drive}) : deux roues motrices indépendantes plus une roue folle assurent à la fois la
propulsion et le pilotage, sans mécanisme de direction dédié. La fiche de projet officielle
résume le concept clé ainsi~:

\begin{conceptbox}{Concept clé du cahier des charges}
	\esibot{} = TurtleBot3 Burger adapté aux ressources de l'ESI. Même architecture
	\ros{}\,2, mêmes paquets (\pkg{slam\_toolbox}, \pkg{Nav2}), mêmes topics --- seuls les
	composants physiques changent~: un ESP32-CAM à la place de l'OpenCR, et un radar
	ultrasonique rotatif (HC-SR04~+~servo) à la place du LiDAR laser.
\end{conceptbox}

Chaque groupe devait réaliser son propre exemplaire d'\esibot{} de façon indépendante sur
une durée de 10 semaines, la comparaison inter-groupes (qualité de cartographie, précision
de navigation, robustesse de détection d'obstacles) constituant la dimension compétitive et
pédagogique de l'exercice.

\section{Comparatif avec le TurtleBot3 Burger (cahier des charges)}
\label{sec:presentation_comparatif}

Le tableau~\ref{tab:tb3_compare_spec} reproduit le comparatif fourni dans la spécification
du projet entre la plate-forme de référence et la version \emph{prévue} d'\esibot{}. Le
chapitre~\ref{ch:archi_globale} présente ensuite l'architecture \emph{réellement
implémentée}, qui diverge de ce tableau sur plusieurs points (notamment le capteur de
cartographie et le contrôleur moteur), analysés en détail au chapitre~\ref{ch:limitations}.

\begin{table}[H]
	\centering
	\caption{Comparatif TurtleBot3 Burger / EsiBot --- tel que spécifié dans le cahier des charges.}
	\label{tab:tb3_compare_spec}
	\renewcommand{\arraystretch}{1.3}
	\begin{tabularx}{\linewidth}{L{3.4cm} L{4.4cm} L{4.8cm}}
		\toprule
		\textbf{Caractéristique} & \textbf{TurtleBot3 Burger}        & \textbf{EsiBot (prévu)}                        \\
		\midrule
		Roues motrices            & 2$\times$ Dynamixel XL430 (servos) & 2$\times$ moteurs DC jaunes + encodeurs        \\
		Roue folle                & 1 ball caster                      & 1 roue folle imprimée 3D                       \\
		Microcontrôleur           & OpenCR 1.0 (Cortex-M7)             & ESP32-CAM (\gls{wifi} intégré + caméra)        \\
		Ordinateur embarqué       & Raspberry Pi 3B+ / 4               & Raspberry\,Pi\,4 (2 ou 4\,Go)                  \\
		Capteur de distance       & LiDAR 360° laser (HLS-LFCD2)        & 1$\times$ HC-SR04 sur servo SG90 (radar rotatif) \\
		Angle de balayage         & 360°                                & 180° (servo) ou 360° (rotation continue)        \\
		Caméra                    & Option (Raspberry Pi Cam)          & ESP32-CAM OV2640 intégrée (streaming \gls{wifi}) \\
		Châssis                   & Pièces ABS moulées ROBOTIS          & Impression 3D PLA conçue par les étudiants     \\
		Framework                 & \ros{} Humble / Iron                & \ros{} Humble (identique)                       \\
		\gls{slam}                & slam\_toolbox / cartographer        & slam\_toolbox (identique)                       \\
		Navigation autonome       & Nav2                                & Nav2 (identique)                                \\
		Open-source matériel      & Oui (STL fournis par ROBOTIS)       & Oui (STL conçus par les étudiants)              \\
		\bottomrule
	\end{tabularx}
\end{table}

L'analogie centrale revendiquée par le cahier des charges est que le radar rotatif
(servo~+~HC-SR04), bien que mécaniquement très différent d'un LiDAR, produit en sortie un
message \texttt{sensor\_msgs/LaserScan} au format strictement identique à celui d'un
LiDAR~: les paquets standards \pkg{slam\_toolbox} et \pkg{Nav2} le consomment donc sans
aucune modification de code. Cette même propriété d'interchangeabilité de capteur est
réutilisée, dans l'implémentation réelle, avec le capteur de profondeur Kinect (voir
chapitre~\ref{ch:slam}).

\section{Liste des livrables du projet}
\label{sec:presentation_livrables}

Le projet définit huit livrables, rappelés au tableau~\ref{tab:livrables}. Ce document est
le \textbf{Livrable n°6}.

\begin{table}[H]
	\centering
	\caption{Livrables attendus par groupe (extrait du cahier des charges).}
	\label{tab:livrables}
	\begin{tabularx}{\linewidth}{C{0.8cm} L{4cm} L{2.6cm} L{5.5cm}}
		\toprule
		\textbf{\#} & \textbf{Livrable} & \textbf{Format} & \textbf{Contenu minimum} \\
		\midrule
		1 & Fichiers \gls{cad} + STL & FreeCAD / STL & Châssis complet : 2 plateaux, entretoises, tourelle servo, roue folle \\
		2 & Schéma électronique & PDF (Fritzing/KiCad) & Câblage ESP32 + L298N + servo + HC-SR04 + RPi4 + alimentation \\
		3 & Workspace \ros{}\,2 & GitHub & Paquets \texttt{esibot\_*} avec launch files, \gls{urdf}, paramètres YAML \\
		4 & Carte \gls{slam} & \texttt{.pgm} + \texttt{.yaml} & Carte d'une zone test de minimum 3\,m $\times$ 3\,m \\
		\rowcolor{EsiGray} 5 & Tableau de bord web & HTML/JS & Carte + vidéo + commandes + angle servo + état batterie \\
		\rowcolor{EsiGray} \textbf{6} & \textbf{Rapport technique} & \textbf{PDF (15--25 p.)} & \textbf{Architecture, schémas, résultats mesurés, analyse limitations, bilan (\emph{ce document})} \\
		7 & Vidéo de démonstration & MP4 (3--5 min) & Démarrage $\rightarrow$ cartographie $\rightarrow$ navigation $\rightarrow$ dashboard \\
		8 & README & Markdown & Installation, lancement, dépendances, architecture \\
		\bottomrule
	\end{tabularx}
\end{table}

\section{Scénario de démonstration imposé}
\label{sec:presentation_scenario}

Le jury place \esibot{} dans une pièce inconnue. L'équipe lance le robot avec une seule
commande. La phase~1 (2~min) consiste en une télé-opération pour construire la carte, qui
apparaît en direct sur le tableau de bord. En phase~2, le jury désigne deux points sur la
carte~; le robot doit alors naviguer seul jusqu'à chaque point en évitant les obstacles, un
obstacle surprise pouvant être ajouté en cours de navigation à titre de bonus. Ce scénario
structure directement le plan de tests du chapitre~\ref{ch:resultats}.
